\documentclass[12pt,a4paper]{article}

\usepackage[margin=1in]{geometry}
\usepackage{amsmath}
\usepackage{booktabs}
\usepackage{listings}
\usepackage{xcolor}
\usepackage{graphicx}
\usepackage{hyperref}
\usepackage{parskip}
\usepackage{microtype}

% Hyperref settings
\hypersetup{
    colorlinks=true,
    linkcolor=blue!60!black,
    urlcolor=blue!60!black,
    citecolor=blue!60!black,
}

% Listings style for pseudocode / VM code
\definecolor{codebg}{RGB}{245,245,245}
\definecolor{codecomment}{RGB}{100,100,100}
\definecolor{codekw}{RGB}{0,80,160}
\definecolor{codestr}{RGB}{160,0,0}

\lstdefinestyle{hackvm}{
    backgroundcolor=\color{codebg},
    basicstyle=\ttfamily\small,
    keywordstyle=\color{codekw}\bfseries,
    commentstyle=\color{codecomment}\itshape,
    stringstyle=\color{codestr},
    breaklines=true,
    frame=single,
    rulecolor=\color{gray!50},
    numbers=left,
    numberstyle=\tiny\color{gray},
    numbersep=8pt,
    tabsize=4,
    showstringspaces=false,
    captionpos=b,
}

\lstdefinestyle{pseudo}{
    backgroundcolor=\color{codebg},
    basicstyle=\ttfamily\small,
    keywordstyle=\color{codekw}\bfseries,
    commentstyle=\color{codecomment}\itshape,
    morekeywords={if, else, return, for, while, function, call, ERROR_LOC},
    breaklines=true,
    frame=single,
    rulecolor=\color{gray!50},
    numbers=left,
    numberstyle=\tiny\color{gray},
    numbersep=8pt,
    tabsize=4,
    showstringspaces=false,
    captionpos=b,
}

\lstset{style=pseudo}

% Section formatting
\usepackage{titlesec}
\titleformat{\section}{\large\bfseries}{}{0em}{\thesection.\quad}
\titleformat{\subsection}{\normalsize\bfseries}{}{0em}{\thesubsection\quad}

% -

\title{%
    \textbf{DA2304 - Introduction to Computer Systems} \\[0.5em]
    \textbf{ Assignment 2}
}
\author{Kiran Kumar P \\
        DA24B008
        }
\date{}

\begin{document}

\maketitle

\hrule

\section*{Extra Credit}

My recent favorite song has been \textit{Stalin Thodaratum Tamil Nadu Vellattum} by \textit{Arivalayam}.  

\noindent
Link: \href{https://youtu.be/cGDoVZmeg5I?si=BldBGriqk_Dvbyse}{May 4: Wait and Watch}

\vspace{1em}
\hrule
\vspace{1em}
% ============================================================
\section{Sanity Checking Logic}
% ============================================================

\subsection{Why dimension validation is necessary?}

For a matrix multiplication $W \times X$, we need \textbf{$W_{\text{cols}} = X_{\text{rows}}$} to hold. If it doesn't, the inner product loop ends up iterating over the wrong number of elements. Since the Hack VM's address space is just a flat array of 16-bit integers with no type system or bounds checking, the loop will happily read whatever values happen to sit next to the matrix in RAM. You get garbage output and no runtime error to tell you something went wrong.

For $B$, we need \textbf{$B_{\text{rows}} = W_{\text{rows}}$}. If this doesn't match, the bias gets added to the wrong rows of the output - again, silent corruption with no way to detect it unless you check beforehand.

\subsection{Validation Pseudocode}

The convention I followed is: write $-1$ to \texttt{ERROR\_LOC} on failure, $1$ on success.

\begin{lstlisting}[style=pseudo, caption={Matrix dimension validation logic}]
function sanity_check(W_rows, W_cols, X_rows, X_cols, B_rows, B_cols):

    // Check: matrix multiplication is defined
    if W_cols != X_rows:
        ERROR_LOC = -1
        return False

    // Check: bias vector has matching row count
    if B_rows != W_rows:
        ERROR_LOC = -1
        return False

    // All dimensions are compatible
    ERROR_LOC = 1
    return True
\end{lstlisting}

For example, if $W$ is $3 \times 4$ and $X$ is $3 \times 2$, then $W_{\text{cols}} = 4 \neq 3 = X_{\text{rows}}$. The check catches this right away and writes $-1$ to \texttt{ERROR\_LOC} before any loop gets a chance to run.

\subsection{Implementation in Hack VM}

In VM code, this is pretty straightforward - push the two dimension arguments, subtract, and branch on the result. Here's the first check:

\begin{lstlisting}[style=hackvm, caption={VM-level dimension mismatch check}]
// if W_cols != X_rows: error
push argument 1    // W_cols
push argument 2    // X_rows
sub
if-goto DIMENSION_ERROR

// ... rest of computation ...

label DIMENSION_ERROR
push constant 1
neg                // push -1
pop static 0       // ERROR_LOC = -1
\end{lstlisting}

The key thing here is that \texttt{if-goto} fires on any nonzero value, so the subtraction result being nonzero means the dimensions don't match. Simple but effective.

% ============================================================
\section{Approach A: MATMUL Analysis}
% ============================================================

\subsection{Architecture}

In Approach A, one function \texttt{MATMUL(W, X, B, ...)} does the entire matrix multiplication by itself. It has a triple-nested loop that iterates over output rows, output columns, and the inner dot-product dimension. No helper functions, no calls - everything lives in one place.

\subsection{Stack Usage Analysis}

Since the whole computation happens inside a single call to \texttt{MATMUL}, the stack frame gets created exactly once. The frame is the return address plus the four saved segment pointers (\texttt{LCL}, \texttt{ARG}, \texttt{THIS}, \texttt{THAT}) - five words total - pushed at call time and restored at return. For an $m \times n$ output matrix with inner dimension $k$, all $m \times n \times k$ accumulate operations happen within this one frame. No extra frames get allocated at any point.

In practice, this means the stack pressure stays low and predictable. SP goes up a bit when the local variables are allocated at function entry, and then it stays more or less flat through the entire triple-nested loop. There are no intermediate \texttt{call}/\texttt{return} pairs creating and tearing down frames, so the segment registers \texttt{LCL}, \texttt{ARG}, \texttt{THIS}, \texttt{THAT} are saved and restored just once. On the Hack VM, each of those save/restore operations is multiple assembly instructions - load the register value into \texttt{D}, write it to the stack, bump \texttt{SP} - so doing it only once actually matters.

\begin{itemize}
    \item \textbf{Frame setups}: 1 total, regardless of matrix size.
    \item \textbf{Segment saves/restores}: 4 each, paid once at call and once at return.
    \item \textbf{Stack growth during loops}: Minimal. Just the local variables allocated at entry.
    \item \textbf{Push/pop traffic inside loops}: Limited to the operands of each arithmetic operation (typically 2 pushes and 1 pop per accumulate step).
\end{itemize}

For large matrices (say $m = n = k = 10$), MATMUL performs 1000 inner iterations with a fixed overhead of one frame creation.

\subsection{Segment Usage Analysis}

MATMUL touches four distinct memory segments, each serving a specific role:

\textbf{Argument segment (ARG):} The matrix base pointers ($W$, $X$, $B$, output, and dimension values) come in as arguments and are accessed via \texttt{ARG} relative addressing. The important thing is that these are read once at function entry and then cached into \texttt{temp} registers, so the \texttt{ARG}-relative fetch cost only shows up during initialization, not inside the hot loop.

\textbf{Local segment (LCL):} Loop counters (row index $i$, column index $j$, inner index $k$) go into local variables. Each local access is an \texttt{LCL}-relative fetch - the VM loads \texttt{LCL}, adds the offset, and dereferences. That's two memory reads per access. I found it helps to put only the less frequently updated values (like the outer loop counter) in locals and keep the hot variables in \texttt{temp}.

\textbf{Temp segment (RAM[5--12]):} This turned out to be the most useful segment for the inner loop. Since \texttt{temp} maps to absolute RAM addresses, accessing \texttt{temp 0} compiles to just \texttt{@5 / D=M} - a single memory read with no pointer indirection. The accumulator and pointer increments belong here. After looking at the generated assembly, the difference between a \texttt{temp} access and a \texttt{local} access was quite noticeable in terms of instruction count. A typical allocation:

\begin{lstlisting}[style=pseudo, caption={Segment allocation for MATMUL}]
temp 0  --> dot product accumulator
temp 1  --> current pointer into row i of W
temp 2  --> current pointer into column j of X
local 0 --> i  (outer loop counter)
local 1 --> j  (middle loop counter)
local 2 --> k  (inner loop counter)
ARG[0]  --> base address of W
ARG[1]  --> base address of X
ARG[2]  --> base address of B
\end{lstlisting}

\textbf{Static/Pointer segments:} Not really used for the core computation. The pointer segment could theoretically hold the \texttt{this}/\texttt{that} base addresses for matrix accesses, but \texttt{temp} works better here since it doesn't interfere with object-oriented conventions.

\subsection{Advantages}

\begin{itemize}
    \item \textbf{Frame overhead is fixed}: One call regardless of $m$, $n$, or $k$.
    \item \textbf{Pointer reuse}: Row and column base pointers increment in \texttt{temp}, no recomputation per iteration.
    \item \textbf{Execution speed}: Fewer instructions per dot-product step compared to Approach B.
\end{itemize}

\subsection{Disadvantages}

\begin{itemize}
    \item \textbf{Large function body}: Hard to unit-test individual rows or columns in isolation.
    \item \textbf{High local variable count}: More locals mean a bigger frame footprint and more complex LCL-relative addressing.
    \item \textbf{Debugging difficulty}: While debugging, I noticed that a pointer arithmetic bug in the inner loop corrupts the entire output matrix at once, and there's no per-call boundary to help narrow down where things went wrong. An off-by-one in the column stride, for example, shifts every output element and it's not obvious which loop level caused it.
\end{itemize}

% ============================================================
\section{Approach B: ROWXCOL Analysis}
% ============================================================

\subsection{Architecture}

In Approach B, the caller handles the outer loops and delegates each dot product to a separate \texttt{ROWXCOL(row\_ptr, col\_ptr, length)} call. This function takes a pointer to a row of $W$, a pointer to a column of $X$, and the length of the vectors, and returns their dot product. The caller iterates over all $(i, j)$ pairs, calls \texttt{ROWXCOL} for each one, adds the bias, and stores the result.

\subsection{Stack Usage Analysis}

Unlike MATMUL, ROWXCOL creates a brand-new stack frame for every single $(i, j)$ output element. Each call pushes 3 arguments onto the stack, then the call mechanism pushes 5 more words (return address, saved \texttt{LCL}, \texttt{ARG}, \texttt{THIS}, \texttt{THAT}), for 8 words of stack growth per call. On return, all 8 get popped.

\begin{lstlisting}[style=pseudo, caption={Stack state during one ROWXCOL call}]
// Before call (caller pushes args):
Stack: [..., row_ptr, col_ptr, length]   // SP += 3

// During call (frame created):
Stack: [..., row_ptr, col_ptr, length,
             ret_addr, saved_LCL, saved_ARG,
             saved_THIS, saved_THAT]     // SP += 5 more

// Inside ROWXCOL (locals allocated):
Stack: [..., frame..., local_acc, local_k]  // SP += num_locals

// On return: entire frame torn down
\end{lstlisting}

For an $m \times n$ output, this whole sequence repeats $m \times n$ times. The cumulative stack traffic scales as $O(m \times n)$ frame operations, each costing roughly 30--50 assembly instructions for the \texttt{call} and a similar count for \texttt{return}. For a $4 \times 4$ output matrix, that's 16 call/return pairs, each saving and restoring 4 segment registers. The segment registers \texttt{LCL}, \texttt{ARG}, \texttt{THIS}, \texttt{THAT} get written to and read from the stack $16 \times 4 = 64$ times just in overhead - before any actual dot product computation even starts.

One thing that became noticeable while stepping through the generated assembly: on an architecture with a dedicated \texttt{CALL} instruction (like x86), the hardware handles register saving in a single cycle. On the Hack VM, it's all done in software, and that cost is very visible when you watch the SP bouncing up and down in the emulator.

\subsection{Segment Usage Analysis}

\textbf{Argument segment (ARG):} ROWXCOL receives three arguments per call: \texttt{row\_ptr}, \texttt{col\_ptr}, and \texttt{length}. These are accessed via ARG-relative addressing (\texttt{push argument 0}, etc.) inside the dot-product loop. Unlike MATMUL which caches its arguments once and reuses them, ROWXCOL has to access \texttt{ARG} fresh on every call. If \texttt{row\_ptr} or \texttt{col\_ptr} is read inside the inner loop body without caching, the ARG-relative fetch cost repeats every $k$ iterations, which adds up fast.

\textbf{Local segment (LCL):} The accumulator and loop counter live as locals inside \texttt{ROWXCOL}. Since \texttt{LCL} is set fresh on every call, these get re-initialized to zero each time. There's no cross-call state in the local segment - which makes correctness easier to reason about but also means you can't reuse any previously computed intermediate values.

\textbf{Temp segment:} ROWXCOL can (and should) cache the argument values into \texttt{temp} at function entry to avoid repeated ARG-relative fetches in the inner loop:

\begin{lstlisting}[style=pseudo, caption={Caching arguments into temp inside ROWXCOL}]
function ROWXCOL 2:
    push argument 0   // row_ptr
    pop temp 0        // cache in temp
    push argument 1   // col_ptr
    pop temp 1        // cache in temp
    push argument 2   // length
    pop local 1       // loop counter
    // inner loop uses temp 0, temp 1
\end{lstlisting}

Without this caching, every loop iteration pays the ARG-relative addressing overhead (load LCL base, add offset, dereference), and it adds up across $k$ iterations and $m \times n$ calls. I noticed this when the inner loop was taking way more instructions than expected - the repeated \texttt{push argument} operations were a big part of it.

\textbf{Cross-call segment pressure:} Because each call creates a new frame, the segment registers \texttt{LCL} and \texttt{ARG} are overwritten on every call and restored on every return. The caller's \texttt{LCL} and \texttt{ARG} get saved to the stack during the call and loaded back on return. This back-and-forth - push before call, pop after return - is a constant source of memory write traffic that MATMUL avoids entirely.

\subsection{Modularity and Reusability}

The real strength of this approach is that \texttt{ROWXCOL} is a self-contained, testable function. It can be:
\begin{itemize}
    \item Unit-tested independently by feeding it known vectors and checking the returned dot product.
    \item Reused across different matrix shapes without any modification.
\end{itemize}

This made debugging much easier - when the output was wrong, I could test \texttt{ROWXCOL} in isolation first and confirm the dot product was correct before looking at the caller's pointer arithmetic.

\subsection{Disadvantages on the Hack VM}

Every \texttt{call} compiles to roughly 30+ assembly instructions (push return address, push 4 saved registers, set ARG, set LCL, jump) and every \texttt{return} to a similar count. For a function called $m \times n$ times, this overhead compounds fast. On a simple accumulator loop that would otherwise be 5--10 instructions per inner iteration, the call overhead can easily dominate actual execution time. After looking at the generated assembly for even a small $3 \times 3$ case, the call/return boilerplate was surprisingly large compared to the actual dot product code.

% ============================================================
\section{Comparative Analysis}
% ============================================================


\begin{table}[h!]
\centering
\caption{MATMUL vs.\ ROWXCOL: Architectural Comparison}

\renewcommand{\arraystretch}{1.2}

\begin{tabular}{@{}p{4cm}p{5.5cm}p{5.5cm}@{}}
\toprule

\textbf{Criterion} &
\textbf{MATMUL (Approach A)} &
\textbf{ROWXCOL (Approach B)} \\

\midrule

Stack frame setups &
1 (for the entire computation) &
$m \times n$ (one per dot product) \\

Argument passing overhead &
Low (pointers loaded once from ARG) &
High (3 arguments re-pushed every call) \\

Segment save/restore &
Once on entry and exit &
Once per \texttt{ROWXCOL} call \\

Execution efficiency &
Better due to fewer redundant instructions &
Worse due to repeated call overhead \\

Memory access pattern &
Pointer arithmetic managed using temp segment &
Pointers recomputed repeatedly for each call \\

Debugging &
Harder to isolate failures &
Easier to debug individual dot-product logic \\

\bottomrule
\end{tabular}

\end{table}

MATMUL optimises for execution efficiency at the cost of structural clarity, while ROWXCOL prioritises correctness and testability at the cost of runtime performance.

% ============================================================
\section{Optimizations}
% ============================================================

\subsection{Optimizations for MATMUL}

\textbf{Precomputing row offsets -}
Instead of recomputing \texttt{base + i * cols} on every outer loop iteration, the row pointer can be kept as an incrementing value in a \texttt{temp} register. Each iteration just adds \texttt{W\_cols} to the pointer, which avoids repeated multiplications. This matters a lot on the Hack VM since multiplication itself has to be done as repeated addition in software - so avoiding a multiply inside the loop saves a bunch of instructions.

\textbf{Reducing unnecessary stack operations -}
Every \texttt{push}/\texttt{pop} pair that can be eliminated saves two memory accesses plus an SP update. While debugging, I noticed a lot of intermediate values were being pushed to the stack only to be popped right back. Keeping those values in \texttt{temp} instead cuts out the round-trip.

\textbf{Loop unrolling -}
If the inner dimension $k$ is known at translation time, the inner loop can be unrolled. For $k=4$:

\begin{lstlisting}[style=hackvm, caption={Manual loop unrolling for k=4 inner dimension}]
// Unrolled dot product: no loop counter needed
push pointer 1        // *row_ptr
push pointer 2        // *col_ptr
add
// ... repeat 3 more times manually
\end{lstlisting}

This gets rid of the loop counter update and the branch check on every iteration. For small fixed dimensions, this can roughly halve the instruction count of the inner loop.

\textbf{Temp register usage strategy -}
The Hack VM's \texttt{temp} segment (RAM[5--12]) is directly addressable without going through a base pointer. Putting frequently written values (accumulators, pointer increments) in \texttt{temp} instead of \texttt{local} eliminates the \texttt{LCL}-relative addressing overhead. In practice, moving the accumulator from \texttt{local 0} to \texttt{temp 0} saved about 4 assembly instructions per inner loop iteration.

\textbf{Hypothetical hardware multiplier -}
The Hack ALU only supports addition and bitwise operations, so any multiplication has to be implemented as repeated addition. If a hardware multiplier were available (say, as a memory-mapped I/O operation), address computations like \texttt{i * cols} could be reduced from $O(\text{cols})$ to $O(1)$ instructions. This would speed up all the array indexing significantly.

\subsection{Optimizations for ROWXCOL}

\textbf{Reducing frame setup overhead -}
The five-word saved frame (return address, LCL, ARG, THIS, THAT) is needed for general VM correctness, but it's partially redundant for a leaf function like \texttt{ROWXCOL} that never calls anything else. A specialized calling convention that only saves the registers actually used (say, just LCL and ARG) would cut the frame from 5 words to 2, reducing setup/teardown cost by about 60\%.

\textbf{Local variable reuse -}
Within \texttt{ROWXCOL}, the accumulator and loop counter should be reused across the loop body rather than re-initialized. Declaring them as locals and updating in-place avoids unnecessary pushes of zero at each iteration start.

\textbf{Minimising repeated argument movement -}
Arguments passed to \texttt{ROWXCOL} are accessed via \texttt{ARG} relative addressing. If the function reads \texttt{argument 0} (row pointer) inside the inner loop body, it pays the \texttt{ARG} relative fetch cost every iteration. Caching the argument into a \texttt{temp} register at function entry avoids this. I actually implemented this in my \texttt{ROWXCOL.vm} and it made a noticeable difference in the instruction count per loop iteration.

\end{document}